
\documentclass{article}
\usepackage{colortbl}
\usepackage{makecell}
\usepackage{multirow}
\usepackage{supertabular}

\begin{document}

\newcounter{utterance}

\centering \large Interaction Transcript for game `guesswhat', experiment `Abs\_Level\_2', episode 0 with NEW\_base\_planning\_clp\_chat.
\vspace{24pt}

{ \footnotesize  \setcounter{utterance}{1}
\setlength{\tabcolsep}{0pt}
\begin{supertabular}{c@{$\;$}|p{.15\linewidth}@{}p{.15\linewidth}p{.15\linewidth}p{.15\linewidth}p{.15\linewidth}p{.15\linewidth}}
    \# & $\;$A & \multicolumn{4}{c}{Game Master} & $\;\:$B\\
    \hline

    \theutterance \stepcounter{utterance}  
    & & \multicolumn{4}{p{0.6\linewidth}}{
        \cellcolor[rgb]{0.9,0.9,0.9}{
            \makecell[{{p{\linewidth}}}]{
                \texttt{\tiny{[P1$\langle$GM]}}
                \texttt{Let's play a guessing game. Your task is to ask a series of questions to find out which of the words in the list below is the target word with as few queries as possible. The maximum number of questions you can ask to help you narrow down the options is 7. Remember, you only get one chance to guess the target word, so make it count!} \\
\\ 
\texttt{The rules of the game are: } \\
\\ 
\texttt{1. Only ask one question per turn. DO NOT ASK MORE THAN ONE QUESTION.} \\
\\ 
\texttt{2. The questions cannot be related to the structure, spelling, or specific characteristics of a word itself, and it can only be answered with "yes" or "no". } \\
\\ 
\texttt{3. Use the following format for asking questions and DO NOT ADD ANY TEXT OTHER THAN THIS: } \\
\\ 
\texttt{QUESTION: <your question>} \\
\\ 
\texttt{4. You have only one opportunity to guess the target word. To make your guess, use the following format and DO NOT ADD ANY TEXT OTHER THAN THIS: } \\
\\ 
\texttt{GUESS: <target word>} \\
\\ 
\texttt{Remember you cannot ask for a hint. Make your question brief. Let’s start!} \\
\\ 
\texttt{List of candidate words: {[}'rage', 'infuriation', 'boredom', 'disappointment', 'earl', 'duke', 'prince', 'princes'{]}} \\
            }
        }
    }
    & \\ \\

    \theutterance \stepcounter{utterance}  
    & \multicolumn{4}{p{0.45\linewidth}}{
        \cellcolor[rgb]{0.95,0.95,0.95}{
            \makecell[{{p{\linewidth}}}]{
                \texttt{\tiny{[P1$|$internal]}}
                \texttt{{[}} \\
\texttt{  \{} \\
\texttt{    "role": "user",} \\
\texttt{    "content": "Let's play a guessing game. Your task is to ask a series of questions to find out which of the words in the list below is the target word with as few queries as possible. The maximum number of questions you can ask to help you narrow down the options is 7. Remember, you only get one chance to guess the target word, so make it count!\textbackslash{}n\textbackslash{}nThe rules of the game are: \textbackslash{}n\textbackslash{}n1. Only ask one question per turn. DO NOT ASK MORE THAN ONE QUESTION.\textbackslash{}n\textbackslash{}n2. The questions cannot be related to the structure, spelling, or specific characteristics of a word itself, and it can only be answered with \textbackslash{}"yes\textbackslash{}" or \textbackslash{}"no\textbackslash{}". \textbackslash{}n\textbackslash{}n3. Use the following format for asking questions and DO NOT ADD ANY TEXT OTHER THAN THIS: \textbackslash{}n\textbackslash{}nQUESTION: <your question>\textbackslash{}n\textbackslash{}n4. You have only one opportunity to guess the target word. To make your guess, use the following format and DO NOT ADD ANY TEXT OTHER THAN THIS: \textbackslash{}n\textbackslash{}nGUESS: <target word>\textbackslash{}n\textbackslash{}nRemember you cannot ask for a hint. Make your question brief. Let\textbackslash{}u2019s start!\textbackslash{}n\textbackslash{}nList of candidate words: {[}'rage', 'infuriation', 'boredom', 'disappointment', 'earl', 'duke', 'prince', 'princes'{]}"} \\
\texttt{  \},} \\
\texttt{  \{} \\
\texttt{    "role": "user",} \\
\texttt{    "content": "You are a skilled planner. Above you see the game history. What would be the best next action? Openly reason about the next step, then produce your next action exactly as the game expects it preceded by the separator 'ACT:' on its own line. If multiple actions are required, write ACT: first and then all actions."} \\
\texttt{  \}} \\
\texttt{{]}} \\
            }
        }
    }
    & & \\ \\

    \theutterance \stepcounter{utterance}  
    & \multicolumn{4}{p{0.45\linewidth}}{
        \cellcolor[rgb]{0.95,0.95,0.95}{
            \makecell[{{p{\linewidth}}}]{
                \texttt{\tiny{[P1$|$internal]}}
                \texttt{ACT:} \\
\texttt{QUESTION: Is the target word a feeling?} \\
            }
        }
    }
    & & \\ \\

    \theutterance \stepcounter{utterance}  
    & \multicolumn{4}{p{0.6\linewidth}}{
        \cellcolor[rgb]{0.9,0.9,0.9}{
            \makecell[{{p{\linewidth}}}]{
                \texttt{\tiny{[P1$\rangle$GM]}}
                \texttt{QUESTION: Is the target word a feeling?} \\
            }
        }
    }
    & & \\ \\

    \theutterance \stepcounter{utterance}  
    & & \multicolumn{4}{p{0.6\linewidth}}{
        \cellcolor[rgb]{0.9,0.9,0.9}{
            \makecell[{{p{\linewidth}}}]{
                \texttt{\tiny{[GM$\rangle$P2]}}
                \texttt{Let’s play a guessing game! Your task is to answer the other player's questions. Based on your knowledge of the word: boredom, respond to the following questions or guesses. Limit your response to only 'yes' or 'no' with no explanation or other words. Never reveal the answer in your response.} \\
\\ 
\texttt{You must reply using the format below and DO NOT ADD ANY TEXT OTHER THAN THIS:} \\
\\ 
\texttt{ANSWER: <some text>} \\
\\ 
\texttt{Target Word: boredom} \\
\\ 
\texttt{QUESTION: Is the target word a feeling?} \\
            }
        }
    }
    & \\ \\

    \theutterance \stepcounter{utterance}  
    & & & \multicolumn{4}{p{0.45\linewidth}}{
        \cellcolor[rgb]{0.95,0.95,0.95}{
            \makecell[{{p{\linewidth}}}]{
                \texttt{\tiny{[P2$|$internal]}}
                \texttt{{[}} \\
\texttt{  \{} \\
\texttt{    "role": "user",} \\
\texttt{    "content": "Let\textbackslash{}u2019s play a guessing game! Your task is to answer the other player's questions. Based on your knowledge of the word: boredom, respond to the following questions or guesses. Limit your response to only 'yes' or 'no' with no explanation or other words. Never reveal the answer in your response.\textbackslash{}n\textbackslash{}nYou must reply using the format below and DO NOT ADD ANY TEXT OTHER THAN THIS:\textbackslash{}n\textbackslash{}nANSWER: <some text>\textbackslash{}n\textbackslash{}nTarget Word: boredom\textbackslash{}n\textbackslash{}nQUESTION: Is the target word a feeling?"} \\
\texttt{  \},} \\
\texttt{  \{} \\
\texttt{    "role": "user",} \\
\texttt{    "content": "You are a skilled planner. Above you see the game history. What would be the best next action? Openly reason about the next step, then produce your next action exactly as the game expects it preceded by the separator 'ACT:' on its own line. If multiple actions are required, write ACT: first and then all actions."} \\
\texttt{  \}} \\
\texttt{{]}} \\
            }
        }
    }
     \\ \\

    \theutterance \stepcounter{utterance}  
    & & & \multicolumn{4}{p{0.45\linewidth}}{
        \cellcolor[rgb]{0.95,0.95,0.95}{
            \makecell[{{p{\linewidth}}}]{
                \texttt{\tiny{[P2$|$internal]}}
                \texttt{ANSWER: yes} \\
\\ 
\texttt{ACT: QUESTION: Is the target word a noun?} \\
            }
        }
    }
     \\ \\

    \theutterance \stepcounter{utterance}  
    & & & \multicolumn{4}{p{0.6\linewidth}}{
        \cellcolor[rgb]{0.9,0.9,0.9}{
            \makecell[{{p{\linewidth}}}]{
                \texttt{\tiny{[GM$\langle$P2]}}
                \texttt{QUESTION: Is the target word a noun?} \\
            }
        }
    }
     \\ \\

    \theutterance \stepcounter{utterance}  
    & & & \multicolumn{2}{p{0.3\linewidth}}{
        \cellcolor[rgb]{0.9,0.9,0.9}{
            \makecell[{{p{\linewidth}}}]{
                \texttt{\tiny{[GM$|$GM]}}
                \texttt{abort game} \\
            }
        }
    }
    & & \\ \\

\end{supertabular}
}

\end{document}
